\documentclass[11pt]{article}
\usepackage[margin=1in]{geometry}
\usepackage{times}
\usepackage{hyperref}
\hypersetup{colorlinks=true, linkcolor=black, urlcolor=blue, citecolor=black}
\title{The Live Demonstration That Went Live: Steorn's Orbo and the Jury That Said No}
\author{Flyxion \\ \small Independent Researcher}
\date{}
\begin{document}
\maketitle

\begin{abstract}
The Irish company Steorn announced in 2006, via a paid advertisement in The Economist, that it had developed a magnetic technology called Orbo capable of producing free, clean energy in violation of the law of conservation of energy, and publicly invited the scientific community to test the claim. This paper reviews the resulting, unusually well-documented sequence of events, an independent jury of scientists convened at Steorn's own invitation, a public demonstration at the Kinetica Museum in London that failed live on camera and was withdrawn mid-event, and a subsequent jury verdict finding no evidence of the claimed energy production, arguing that this case is valuable precisely because it is one of the few in this genre where the claimant invited rigorous independent scrutiny at the outset and the resulting negative verdict was unambiguous, public, and never seriously disputed on its own procedural terms even by the company itself.
\end{abstract}

\section{Introduction}

Steorn, a Dublin-based technology company, took out a full-page advertisement in The Economist in August 2006 announcing that it had developed a technology, later named Orbo, using arrangements of permanent magnets said to produce a net energy gain in violation of the conservation of energy, and explicitly invited scientists to apply to an independent jury that would evaluate the technology under controlled conditions, an unusually direct and public invitation to exactly the kind of scrutiny most devices in this series conspicuously avoid.

\section{The Jury Process}

Steorn convened a jury of physicists and engineers, drawn from applicants responding to its public call, who reviewed technical documentation and test data Steorn provided over an extended evaluation period beginning in 2007. In 2009, the jury issued its verdict: it had found no evidence that Orbo produced energy, and it specifically noted that Steorn had not, over more than two years of evaluation, provided a test that unambiguously and reproducibly demonstrated the claimed effect under conditions the jury considered adequately controlled, a negative and specific finding rather than an inconclusive or deadlocked one.

\section{The 2007 Public Demonstration Failure}

Independent of the jury process, Steorn staged a live public demonstration of an Orbo device at the Kinetica Museum in London in July 2007, intended to run continuously and be viewable in person and via live webcast; the device stopped working during the demonstration itself, and Steorn withdrew it citing overheating from the exhibit's stage lighting, an explanation that did not prevent the failure from being witnessed live by the attending public and press, producing a widely covered and unambiguous public record of the specific demonstration not working as claimed, independent of any dispute over the jury's later, separately conducted technical evaluation.

\section{Why This Case Is Comparatively Clean, Evidentially}

Unlike several of the claims surveyed elsewhere in this series, Steorn's own conduct removed most of the usual ambiguity: the company publicly invited independent evaluation rather than avoiding it, submitted to a jury process of its own design and selection, and the resulting negative finding was specific and was not accompanied by a persuasive company rebuttal disputing the jury's competence or methodology in the way over-unity claim promoters more commonly respond to negative results, though Steorn did continue limited further promotion of Orbo in subsequent years, including further demonstrations of comparable ambiguity, before winding the company down amid financial difficulty in 2013. The initial willingness to invite scrutiny, followed by a clear negative result the company did not seriously contest on procedural grounds, makes this a useful contrast case against the pattern, common elsewhere in this series, of claims sustained primarily by never permitting the test that would settle them.

\section{Why the Underlying Physics Objection Still Applies}

Independent of the jury and demonstration history, Orbo's claimed mechanism, permanent magnet interactions producing net energy, is subject to the same conservation-of-energy objection addressed for the Yildiz Motor, Perendev Motor, and other permanent-magnet over-unity claims elsewhere in this series: a static or slowly moving permanent magnet field does no net work over a closed interaction cycle, and no published Orbo technical description has specified a mechanism that avoids this constraint rather than merely asserting that it does.

\section{The Entropic Gravity Framing}

Orbo makes no gravitational claim, concerning conservation of energy via magnetic interaction alone, and is included in this series as a comparative case study in the over-unity genre, notable primarily for the unusually direct and public nature of its own invited scrutiny and the correspondingly unambiguous negative result that scrutiny produced.

\section{Objections}

Steorn's founder maintained in public statements following the jury's verdict that the technology remained real and that the jury's specific test conditions had not been appropriate for demonstrating the effect, a position the company did not, however, substantiate with a subsequent independently verified demonstration meeting an evidentiary standard the jury or other independent observers accepted as adequate before the company's eventual closure.

\section{Conclusion}

Steorn's Orbo case is distinguished from most others in this series by the company's own initial willingness to invite independent, public scrutiny, which produced a comparatively clean and specific negative result, both from a formally convened scientific jury and from a live public demonstration failure, without the company mounting a substantiated rebuttal before winding down operations, leaving the underlying permanent-magnet conservation-of-energy objection as unaddressed here as in the other over-unity claims surveyed in this series.

\end{document}
